\section{Adjunctions, equivalences, and preservation}\label{sec:adjunctions}

\subsection{Two equivalent descriptions of an adjunction}

Let $F:C\to D$ and $G:D\to C$ be fixed presented functors.  A
\emph{displayed hom-adjunction} consists of primitive recursive operations
\[
 \Phi_{a,b}:\Hom_D(Fa,b)\longrightarrow\Hom_C(a,Gb),\qquad
 \Psi_{a,b}:\Hom_C(a,Gb)\longrightarrow\Hom_D(Fa,b),
\]
respecting equality, inverse modulo equality, and natural in the two
variables.  For $u:a'\to a$, $v:b\to b'$, and $f:Fa\to b$, naturality means
\[
 \Phi_{a',b'}(v f F(u))=G(v)\Phi_{a,b}(f)u.
\]
The displayed hom-sets are notation for typed codes modulo the chosen
arithmetic relation.  No set of all maps between them is postulated.

\begin{theorem}[Hom-adjunction and triangle identities]\label{thm:adjunction}
A displayed hom-adjunction is equivalently given by natural transformations
\[
 \eta:1_C\Longrightarrow GF,\qquad
 \varepsilon:FG\Longrightarrow1_D
\]
with primitive recursive component operations, satisfying
\[
 \varepsilon_{Fa}F(\eta_a)=\id_{Fa},\qquad
 G(\varepsilon_b)\eta_{Gb}=\id_{Gb}.
\]
The conversion in each direction is primitive recursive and its correctness
uses only finite intuitionistic equational reasoning.
\end{theorem}
\begin{proof}
Starting with the unit and counit, define
\[
 \Phi(f)=G(f)\eta_a,\qquad \Psi(g)=\varepsilon_bF(g).
\]
Both terms have the stated types and respect equality.  Naturality of the
counit at $f:Fa\to b$, followed by the first triangle, gives
\[
 \Psi\Phi(f)
 =\varepsilon_bFG(f)F(\eta_a)
 =f\varepsilon_{Fa}F(\eta_a)=f.
\]
Naturality of the unit at $g:a\to Gb$, followed by the second triangle,
gives
\[
 \Phi\Psi(g)
 =G(\varepsilon_b)GF(g)\eta_a
 =G(\varepsilon_b)\eta_{Gb}g=g.
\]
For naturality in the source, use $GF(u)\eta_{a'}=\eta_a u$:
\[
 \Phi(fF(u))=G(f)GF(u)\eta_{a'}=G(f)\eta_a u=\Phi(f)u.
\]
Naturality in the target is
$\Phi(vf)=G(v)G(f)\eta_a=G(v)\Phi(f)$.  Combining these proves the stated
bivariable law.

Conversely, let $\Phi,\Psi$ be inverse natural correspondences and set
\[
 \eta_a=\Phi_{a,Fa}(\id_{Fa}),\qquad
 \varepsilon_b=\Psi_{Gb,b}(\id_{Gb}).
\]
Target naturality immediately gives
$\Phi(f)=G(f)\eta_a$.  The inverse correspondence is also natural: for
example, applying $\Phi$ to $\Psi(g)F(u)$ gives $gu$, so injectivity of
$\Phi$ gives $\Psi(gu)=\Psi(g)F(u)$.  The same argument in the target gives
$\Psi(G(v)g)=v\Psi(g)$.  Taking $g=\id_{Gb}$ in the source formula proves
$\Psi(g)=\varepsilon_bF(g)$ for arbitrary $g:a\to Gb$.

The inverse equations at the relevant identities now read
\[
 \varepsilon_{Fa}F(\eta_a)
 =\Psi(\eta_a)=\id_{Fa},\qquad
 G(\varepsilon_b)\eta_{Gb}
 =\Phi(\varepsilon_b)=\id_{Gb}.
\]
For $u:a\to a'$, source naturality and the formula for $\Phi$ yield
$\eta_{a'}u=\Phi(F(u))=GF(u)\eta_a$.  For $v:b\to b'$, target naturality
of $\Psi$ and its formula give
$v\varepsilon_b=\Psi(G(v))=\varepsilon_{b'}FG(v)$.  These are precisely
the naturality equations for the unit and counit.  The two conversion
procedures recover the original data modulo their appropriate equalities.
\end{proof}

This proof explains why existence claims need care.  A family of statements
$\forall a,b\ \exists$ a bijection between hom-collections does not itself
supply the primitive recursive operations $\Phi,\Psi$.  Conversely, once
those operations are supplied, no choice principle is needed in the theorem.

\begin{proposition}[Uniqueness of a specified right adjoint]
For fixed $F$, any two displayed right adjoints $G,G'$ are naturally
isomorphic by explicit component operations.
\end{proposition}
\begin{proof}
Use primes for the second correspondence and counit.  Define
\[
 u_b=\Phi'_{Gb,b}(\varepsilon_b):Gb\to G'b,
 \qquad
 v_b=\Phi_{G'b,b}(\varepsilon'_b):G'b\to Gb.
\]
Then $\Psi'(u_b)=\varepsilon_b$ and $\Psi(v_b)=\varepsilon'_b$.
Naturality of $\Psi$ in its source gives
\[
 \Psi(v_bu_b)=\Psi(v_b)F(u_b)
 =\varepsilon'_bF(u_b)=\Psi'(u_b)=\varepsilon_b
 =\Psi(\id_{Gb}).
\]
Injectivity of $\Psi$ proves $v_bu_b=\id_{Gb}$; the other composite follows
by interchanging primes.  For $h:b\to b'$, applying $\Psi'$ to
$G'(h)u_b$ gives $h\varepsilon_b$.  Applying it to $u_{b'}G(h)$ gives
$\varepsilon_{b'}FG(h)$, which is the same arrow by naturality of
$\varepsilon$.  Thus $G'(h)u_b=u_{b'}G(h)$, proving naturality.
\end{proof}

\subsection{Adjoints preserve supplied universal constructions}

\begin{theorem}[Preservation of limits and colimits]\label{thm:adjoint-preserve}
A displayed right adjoint preserves every supplied fixed finite limit, with
an explicitly constructed primitive recursive mediating operation.  A
displayed left adjoint preserves every supplied fixed finite colimit.
\end{theorem}
\begin{proof}
Suppose $(p_j:l\to H_j)$ is a chosen limit of a fixed finite diagram
$H:J\to D$.  The arrows $G(p_j):Gl\to GH_j$ form a cone by functoriality.
Given a cone $c_j:x\to GH_j$, transpose its components to
$\Psi(c_j):Fx\to H_j$.  For $h:i\to j$ in $J$, naturality gives
\[
 H(h)\Psi(c_i)=\Psi(GH(h)c_i)=\Psi(c_j),
\]
so this is a cone to $H$.  Let $t:Fx\to l$ be its supplied limiting
mediator and put $m=\Phi(t):x\to Gl$.  Then
\[
 G(p_j)m=\Phi(p_jt)=\Phi\Psi(c_j)=c_j.
\]
If $v:x\to Gl$ has the same projections, source and target naturality of
the inverse correspondence give
$p_j\Psi(v)=\Psi(G(p_j)v)=\Psi(c_j)$ for every $j$.
Limit uniqueness gives $\Psi(v)=t$, hence $v=\Phi(t)=m$.

This is an explicit formula for the new mediator: transpose the finite
input tuple, apply the given mediator, and transpose back.  All three stages
are primitive recursive.  Applying the proof in the opposite categories
proves the colimit assertion.  The same algebra applies to an infinite
indexed cone family only when the family and its universal operation have
already been presented with the requisite uniform primitive recursive
parameters; no general infinite-limit existence assertion is being made.
\end{proof}

\subsection{A constructive equivalence criterion}

A functor $F:C\to D$ is \emph{faithful} if $F(f)=F(g)$ implies $f=g$ for
parallel arrows.  Its \emph{split fullness data} are a fixed primitive
recursive operation
\[
 L_{a,a'}:\Hom_D(Fa,Fa')\longrightarrow\Hom_C(a,a')
 \quad\hbox{with}\quad F(L(h))=h.
\]
Together these say that every hom-map has a displayed inverse.  No test for
arrow equality is required.  \emph{Split essential surjectivity} means
specified primitive recursive data
\[
 s(b)\in C,\quad e_b:Fs(b)\to b,\quad d_b:b\to Fs(b),
 \qquad e_bd_b=\id_b,\quad d_be_b=\id_{Fs(b)}.
\]

\begin{theorem}[Split equivalence criterion]\label{thm:split-equivalence}
Faithfulness, split fullness, and split essential surjectivity construct a
primitive recursively presented quasi-inverse $G:D\to C$ and natural
isomorphisms $1_C\cong GF$, $FG\cong1_D$.  They may be chosen to satisfy
the triangle identities, so the result is an adjoint equivalence.
\end{theorem}
\begin{proof}
Put
\[
 G(b)=s(b),\qquad G(v:b\to b')=L_{s(b),s(b')}(d_{b'}ve_b).
\]
By split fullness,
$FG(v)=d_{b'}ve_b$.  Applying this to an identity and using faithfulness
shows $G(\id_b)=\id_{Gb}$.  For $v:b\to b'$ and $w:b'\to b''$,
\[
 F(G(w)G(v))
 =d_{b''}we_{b'}d_{b'}ve_b
 =d_{b''}wve_b=FG(wv).
\]
Faithfulness proves composition preservation.  Equality preservation of
$G$ follows by the same use of faithfulness and congruence.

The maps $e_b:FGb\to b$ are natural, since
$e_{b'}FG(v)=e_{b'}d_{b'}ve_b=ve_b$.  Their displayed inverses are $d_b$.
Define
\[
 \eta_a=L_{a,s(Fa)}(d_{Fa}),\qquad
 \theta_a=L_{s(Fa),a}(e_{Fa}).
\]
Their images under $F$ are inverse, so faithfulness proves
$\theta_a\eta_a=\id_a$ and $\eta_a\theta_a=\id_{GFa}$.  For
$u:a\to a'$, the two sides of $GF(u)\eta_a=\eta_{a'}u$ have images
\[
 d_{Fa'}F(u)e_{Fa}d_{Fa}=d_{Fa'}F(u),
\]
and hence are equal.  The inverse component families are natural as well:
multiplying the counit naturality equation by $d_{b'}$ and $d_b$ gives
$FG(v)d_b=d_{b'}v$, and multiplying the unit equation by $\theta_{a'}$
and $\theta_a$ gives $u\theta_a=\theta_{a'}GF(u)$.
Thus $\eta$ and $e$ are natural isomorphisms with displayed natural inverses.

The first triangle is
$e_{Fa}F(\eta_a)=e_{Fa}d_{Fa}=\id_{Fa}$.
For the second, apply $F$ to $G(e_b)\eta_{Gb}$ and calculate
\[
 F(G(e_b)\eta_{Gb})
 =d_be_be_{FGb}d_{FGb}=d_be_b=\id_{FGb}.
\]
Faithfulness yields $G(e_b)\eta_{Gb}=\id_{Gb}$.  All constructed operations
are finite compositions of the supplied primitive recursive operations.
\end{proof}

Conversely, for an adjoint equivalence with inverse components $\eta_a^{-1}$,
the formula $L(h)=\eta_{a'}^{-1}G(h)\eta_a$ supplies split fullness.  Indeed,
the first triangle says $F(\eta_a)=e_{Fa}^{-1}$; naturality of $e$ gives
$F(L(h))=e_{Fa'}FG(h)e_{Fa}^{-1}=h$.  If $F(f)=F(g)$, then naturality of
$\eta$ gives
$f=\eta_{a'}^{-1}GF(f)\eta_a
 =\eta_{a'}^{-1}GF(g)\eta_a=g$.
The quasi-inverse and counit supply split essential surjectivity.

The adjective ``split'' is substantive.  A mere existential preimage for
each object or arrow is not automatically a primitive recursive selector.
The usual informal phrase ``choose a preimage for every object'' is replaced
here by the displayed data that make that step legitimate.

\section{Monads and their canonical categories}\label{sec:monads}

\subsection{The monad of an adjunction}

A \emph{presented monad} consists of a fixed presented endofunctor $T:C\to C$
and natural transformations $\eta:1_C\Rightarrow T$ and
$\mu:T^2\Rightarrow T$ with primitive recursive components and laws
\[
 \mu_aT(\eta_a)=\id_{Ta}=\mu_a\eta_{Ta},\qquad
 \mu_aT(\mu_a)=\mu_a\mu_{Ta}.
\]
All transformation equalities are componentwise hom-setoid equalities.

\begin{theorem}[Every displayed adjunction induces a monad]
For $F\dashv G$ with unit $\eta$ and counit $\varepsilon$, the operations
\[
 T=GF,\qquad \mu_a=G(\varepsilon_{Fa})
\]
together with $\eta$ form a presented monad.
\end{theorem}
\begin{proof}
Functoriality of $T$ follows by composition.  Naturality of $\mu$ is the
image under $G$ of naturality of $\varepsilon$ at arrows $F(f)$.
The first unit equation is the image under $G$ of the first triangle:
\[
 \mu_aT(\eta_a)=G(\varepsilon_{Fa}F(\eta_a))=\id_{GFa}.
\]
The other is the second triangle at $Fa$:
$\mu_a\eta_{Ta}=G(\varepsilon_{Fa})\eta_{GFa}=\id_{GFa}$.
For associativity, naturality of $\varepsilon$ at
$\varepsilon_{Fa}:FGFa\to Fa$ states
\[
 \varepsilon_{Fa}FG(\varepsilon_{Fa})
 =\varepsilon_{Fa}\varepsilon_{FGFa}.
\]
Applying $G$ gives exactly $\mu_aT(\mu_a)=\mu_a\mu_{Ta}$.
\end{proof}

\subsection{Kleisli category and its adjunction}

\begin{theorem}[Kleisli construction]\label{thm:kleisli}
A presented monad $T$ determines a presented category $C_T$ with the same
object codes as $C$, arrows $a\to b$ represented by arrows $a\to Tb$ in
$C$, identity $\eta_a$, and composition
\[
 g\mathbin{\star} f=\mu_cT(g)f
 \quad(f:a\to Tb,\ g:b\to Tc).
\]
There is a displayed adjunction $F_T\dashv U_T$ that induces the original
monad.
\end{theorem}
\begin{proof}
The arrow predicate is $\Hom_C(a,Tb,f)$ and equality is the corresponding
arithmetic equality formula in $C$.  All operations have primitive recursive
code definitions.  Typing and congruence follow from those in $C$ and $T$.
The two identity equations are
\[
 \eta_b\star f=\mu_bT(\eta_b)f=f,
 \qquad
 f\star\eta_a=\mu_bT(f)\eta_a
 =\mu_b\eta_{Tb}f=f.
\]
For $h:c\to Td$, naturality of $\mu$ at $h$ gives
$T(h)\mu_c=\mu_{Td}T^2(h)$.  Thus
\begin{align*}
 h\star(g\star f)
 &=\mu_dT(h)\mu_cT(g)f\\
 &=\mu_d\mu_{Td}T^2(h)T(g)f\\
 &=\mu_dT(\mu_d)T^2(h)T(g)f\\
 &=\mu_dT(\mu_dT(h)g)f=(h\star g)\star f.
\end{align*}

Define $F_T(a)=a$, $F_T(f:a\to b)=\eta_bf$, and
$U_T(a)=Ta$, $U_T(f:a\to Tb)=\mu_bT(f)$.  Their identity laws are the
monad unit laws.  For $f:a\to b$, $g:b\to c$,
\[
 F_T(g)\star F_T(f)
 =\mu_cT(\eta_c)T(g)\eta_bf
 =T(g)\eta_bf=\eta_cgf=F_T(gf).
\]
For Kleisli arrows $f:a\to Tb$, $g:b\to Tc$, associativity and naturality
of $\mu$ give
\begin{align*}
 U_T(g\star f)
 &=\mu_cT(\mu_c)T^2(g)T(f)\\
 &=\mu_c\mu_{Tc}T^2(g)T(f)\\
 &=\mu_cT(g)\mu_bT(f)=U_T(g)U_T(f).
\end{align*}

The two sides of the hom correspondence
\[
 \Hom_{C_T}(F_Ta,b)\longleftrightarrow\Hom_C(a,U_Tb)
\]
are represented by the same codes and the same equality, so take both maps
to be the identity.  They are natural: precomposition by $F_T(u)$ in the
Kleisli category sends $f$ to
$\mu_bT(f)\eta_a u=fu$, and postcomposition by a Kleisli arrow $v$ sends
$f$ to $\mu_cT(v)f=U_T(v)f$.  The unit is the original $\eta$; the counit
at $b$ is the Kleisli arrow $Tb\to b$ represented by $\id_{Tb}$.
Its image under $U_T$ is $\mu_b$.  Moreover
$U_TF_T(f)=\mu_bT(\eta_b)T(f)=T(f)$, so the induced monad is the
original one modulo the displayed hom equalities.
\end{proof}

\subsection{Eilenberg--Moore category and its adjunction}

A $T$-algebra is a pair $(a,h)$ with $h:Ta\to a$ satisfying
\[
 h\eta_a=\id_a,\qquad h\mu_a=hT(h).
\]
A morphism $(a,h)\to(b,k)$ is an arrow $f:a\to b$ satisfying
$f h=kT(f)$.  These are arithmetical predicates on number codes, not
new comprehension axioms.

\begin{theorem}[Eilenberg--Moore construction]\label{thm:em}
The displayed $T$-algebras and their morphisms form a presented category
$C^T$.  There is a displayed adjunction $F^T\dashv U^T$, where $U^T$ forgets
the algebra operation and
\[
 F^T(a)=(Ta,\mu_a),\qquad F^T(f)=T(f).
\]
This adjunction induces the original monad.
\end{theorem}
\begin{proof}
The identity arrow is an algebra morphism because $\id_a h=hT(\id_a)$.
For algebra morphisms $f:(a,h)\to(b,k)$ and $g:(b,k)\to(c,l)$,
\[
 (gf)h=g(kT(f))=lT(g)T(f)=lT(gf).
\]
Associativity and congruence are inherited from $C$.  Thus the predicates
and operations define a presented category.

The monad laws make $(Ta,\mu_a)$ an algebra.  Naturality of $\mu$ makes
$T(f)$ an algebra morphism, so $F^T$ is a functor; its other laws and the
functor laws for $U^T$ are immediate from the corresponding supplied laws of
$T$ and $C$.

For an algebra $(b,k)$ define the correspondence
\[
 \Hom_{C^T}((Ta,\mu_a),(b,k))\longleftrightarrow\Hom_C(a,b)
\]
by $f\mapsto f\eta_a$ and $g\mapsto kT(g)$.  The second expression is
an algebra morphism because
\[
 kT(g)\mu_a=k\mu_bT^2(g)
 =kT(k)T^2(g)=kT(kT(g)).
\]
Its composite with the first map sends $g$ to
$kT(g)\eta_a=k\eta_bg=g$.
Conversely, if $f$ is an algebra morphism, the other composite is
\[
 kT(f\eta_a)=kT(f)T(\eta_a)
 =f\mu_aT(\eta_a)=f.
\]
For $u:a'\to a$, precomposition by $F^T(u)=T(u)$ changes
$f\eta_a$ to $fT(u)\eta_{a'}=f\eta_a u$.
For an algebra morphism $v:(b,k)\to(b',k')$, postcomposition changes it to
$vf\eta_a$, exactly as required.  Hence the correspondence is natural.
The unit is $\eta$ and the counit at $(b,k)$ is the algebra morphism
$k:(Tb,\mu_b)\to(b,k)$.  Its underlying arrow at the free algebra is
$\mu_a$, proving the final assertion.
\end{proof}

\subsection{Algebra limits without a preservation assumption}

\begin{theorem}[The forgetful functor creates supplied finite limits]
Let $H:J\to C^T$ be a fixed finite diagram.  If its underlying diagram in
$C$ has a supplied limit, it has a canonically supplied algebra structure
making it a limit in $C^T$.  The construction is primitive recursive and
does not require $T$ to preserve that limit.
\end{theorem}
\begin{proof}
Write $H_j=(a_j,h_j)$ and let $(p_j:l\to a_j)$ be the underlying limiting
cone.  The arrows $h_jT(p_j):Tl\to a_j$ form a cone.  Indeed, for an
algebra morphism $H(u):(a_i,h_i)\to(a_j,h_j)$,
\[
 H(u)h_iT(p_i)=h_jT(H(u))T(p_i)=h_jT(p_j).
\]
The supplied limiting operation therefore gives a unique $h:Tl\to l$
with $p_jh=h_jT(p_j)$.

To prove the unit law, compose with every projection:
\[
 p_jh\eta_l=h_jT(p_j)\eta_l
 =h_j\eta_{a_j}p_j=p_j.
\]
Limit uniqueness gives $h\eta_l=\id_l$.  To prove associativity, the two
candidate arrows $TTl\to l$ have the following projections:
\begin{align*}
 p_jh\mu_l
 &=h_jT(p_j)\mu_l
  =h_j\mu_{a_j}T^2(p_j)
  =h_jT(h_j)T^2(p_j),\\
 p_jhT(h)
 &=h_jT(p_j)T(h)
  =h_jT(p_jh)
  =h_jT(h_j)T^2(p_j).
\end{align*}
Again uniqueness proves $h\mu_l=hT(h)$.  Thus $(l,h)$ is an algebra and
its projections are algebra morphisms by their defining equations.

Finally, a cone of algebra morphisms $c_j:(x,k)\to(a_j,h_j)$ has a unique
underlying mediator $u:x\to l$.  For every $j$,
\[
 p_ju k=c_jk=h_jT(c_j)
 =h_jT(p_j)T(u)=p_jhT(u).
\]
Consequently $uk=hT(u)$ by underlying limit uniqueness, so $u$ is an
algebra morphism.  Its uniqueness in $C^T$ follows from uniqueness in $C$.
Every construction used the given finite tuple of projections, $T$, the
algebra operations, and the supplied underlying mediator; all are primitive
recursive.
\end{proof}

These results supply a substantial structural foundation: adjunctions both
produce monads and arise canonically from them, while algebraic structure
can be transported across existing universal constructions.  They do not
assert a general monadicity theorem, an adjoint functor theorem, or the
existence of arbitrary infinite colimits.  Such statements have additional
hypotheses whose arithmetic content must be separately exhibited.
